\documentclass[aspectratio=169,11pt]{beamer}

% ============================================
% PACKAGES
% ============================================
\usepackage[utf8]{inputenc}
\usepackage[T1]{fontenc}
\usepackage[ngerman]{babel}
\usepackage{lmodern}
\usepackage{tcolorbox}
\tcbuselibrary{skins, hooks}
\usepackage{listings}
\usepackage{xcolor}
\usepackage{fontawesome5}

% Modern Font (if installed, otherwise falls back smoothly)
\IfFileExists{FiraSans.sty}{\usepackage[sfdefault]{FiraSans}}{}
\renewcommand{\familydefault}{\sfdefault}

% ============================================
% COLORS & DESIGN SYSTEM
% ============================================
\definecolor{BgColor}{HTML}{FAFAFA}
\definecolor{Primary}{HTML}{2B3A42}
\definecolor{Accent}{HTML}{E84A5F}
\definecolor{Secondary}{HTML}{3F5765}
\definecolor{CodeBg}{HTML}{F0F0F0}
\definecolor{Success}{HTML}{28A745}

\setbeamercolor{background canvas}{bg=BgColor}
\setbeamercolor{title}{fg=Primary}
\setbeamercolor{frametitle}{bg=Primary, fg=white}
\setbeamercolor{structure}{fg=Accent}
\setbeamercolor{normal text}{fg=Primary}
\setbeamercolor{itemize item}{fg=Accent}
\setbeamercolor{itemize subitem}{fg=Secondary}

\setbeamertemplate{navigation symbols}{}
\setbeamertemplate{footline}{%
  \leavevmode%
  \hbox{%
  \begin{beamercolorbox}[wd=.333333\paperwidth,ht=2.25ex,dp=1ex,center]{author in head/foot}%
    \usebeamerfont{author in head/foot}\tiny Falko Himmel
  \end{beamercolorbox}%
  \begin{beamercolorbox}[wd=.333333\paperwidth,ht=2.25ex,dp=1ex,center]{title in head/foot}%
    \usebeamerfont{title in head/foot}\tiny Grundlagen der Praktischen Informatik
  \end{beamercolorbox}%
  \begin{beamercolorbox}[wd=.333333\paperwidth,ht=2.25ex,dp=1ex,right]{date in head/foot}%
    \usebeamerfont{date in head/foot}\insertframenumber{} / \inserttotalframenumber\hspace*{2ex} 
  \end{beamercolorbox}}%
  \vskip0pt%
}

% ============================================
% CUSTOM BOXES
% ============================================
\tcbset{
    modernbox/.style={
        enhanced,
        boxrule=0pt,
        frame hidden,
        colback=white,
        coltitle=white,
        colbacktitle=Secondary,
        fonttitle=\bfseries,
        drop lifted shadow,
        arc=4pt,
        toptitle=2mm,
        bottomtitle=2mm,
        left=2mm,
        right=2mm
    },
    alertbox/.style={
        modernbox,
        colbacktitle=Accent
    },
    successbox/.style={
        modernbox,
        colbacktitle=Success
    }
}

% ============================================
% CODE LISTINGS
% ============================================
\lstdefinelanguage{Haskell}{
  keywords={class, data, deriving, do, else, if, in, import, let, module, of, then, type, where, instance},
  keywordstyle=\color{Accent}\bfseries,
  ndkeywords={print, main, return},
  ndkeywordstyle=\color{Secondary}\bfseries,
  identifierstyle=\color{Primary},
  sensitive=true,
  comment=[l]{--},
  commentstyle=\color{gray}\ttfamily\itshape,
  stringstyle=\color{teal}\ttfamily,
  morestring=[b]"
}

\lstset{
    language=Haskell,
    backgroundcolor=\color{CodeBg},
    basicstyle=\ttfamily\small,
    breaklines=true,
    frame=single,
    rulecolor=\color{CodeBg},
    frameround=tttt,
    numbers=left,
    numberstyle=\tiny\color{gray},
    xleftmargin=15pt,
    tabsize=4,
    showstringspaces=false
}

% ============================================
% TITLE PAGE SETTINGS
% ============================================
\title{\textbf{Tutorium 2}}
\subtitle{Grundlagen Haskell}
\institute{Grundlagen der Praktischen Informatik}
\author{}
\date{\today}

\begin{document}

\begin{frame}[plain]
    \titlepage
\end{frame}

\begin{frame}{Inhaltsverzeichnis}
    \tableofcontents
\end{frame}

% ============================================
% THEORIE
% ============================================
\section{Theorie: Funktionen \& Typen}

\begin{frame}[fragile]{Funktionen in Haskell}
    Funktionen in Haskell verhalten sich ähnlich wie mathematische Funktionen. Sie bestehen aus einem \textbf{Bezeichner}, \textbf{Parametern} und einem \textbf{Ausdruck}.
    
    \begin{tcolorbox}[modernbox, title=Beispiel: Addition]
\begin{lstlisting}
-- Definition ohne explizite Typen
add x y = x + y
\end{lstlisting}
    \end{tcolorbox}
    
    \begin{itemize}
        \item \textbf{Bezeichner:} \texttt{add}
        \item \textbf{Parameter:} \texttt{x} und \texttt{y}
        \item \textbf{Ausdruck:} \texttt{x + y}
    \end{itemize}
\end{frame}

\begin{frame}[fragile]{Typsignaturen}
    In Haskell kann (und sollte) man optional eine \textbf{Typsignatur} angeben.
    
    \begin{tcolorbox}[modernbox, title=Funktion mit Typsignatur]
\begin{lstlisting}
-- func' :: Parameter 1 -> Parameter 2 -> Parameter 3 -> Rueckgabetyp
func' :: Integer -> Integer -> String -> Integer
func' x y z = x + y
\end{lstlisting}
    \end{tcolorbox}
    
    \begin{itemize}
        \item \texttt{::} bedeutet ``hat den Typ''.
        \item Die Typen der Parameter werden nacheinander mit \texttt{->} getrennt.
        \item Der \textbf{letzte Typ} in der Kette ist immer der \textbf{Rückgabetyp}.
    \end{itemize}
\end{frame}

\begin{frame}[fragile]{Rekursion Basics}
    Ein rekursiver Algorithmus benötigt immer einen \textbf{Basisfall}, damit er terminieren kann. Ansonsten entsteht eine Endlosschleife!
    
    \begin{tcolorbox}[alertbox, title=Beispiel: Fakultät]
\begin{lstlisting}
-- 5! = 5 * 4! = 5 * (4 * 3!) ...
fak :: Integer -> Integer
fak 0 = 1                   -- Basisfall
fak x = x * fak (x-1)       -- Rekursiver Aufruf
\end{lstlisting}
    \end{tcolorbox}
\end{frame}

% ============================================
% LIVE DEMO
% ============================================
\section{Praxis}

\begin{frame}[plain]
    \begin{center}
        \Huge \textbf{Live Demo} \\
        \vspace{0.5cm}
        \large \faLaptopCode~ \texttt{haskell.hs} \\
        \vspace{0.3cm}
        \normalsize Wir programmieren Pattern Matching und verschiedene Rekursionsarten zusammen im Code!
    \end{center}
    
    \vspace{0.5cm}
    \textbf{Aufgaben:}
    \begin{enumerate}
        \item Wir wollen die Zahlen 1 - 3 als Wort ausgeben (,,eins``, ,,zwei``), sonst einen Fehler.
        \item Fakultät mit Akkumulator (Endrekursion) definieren.
        \item Welche Rekursionsarten liegen vor?
    \end{enumerate}
\end{frame}

\end{document}
